% ---------------------------------------------------------------
% Revision text for Review 2: repetitions, warm-up, run-to-run variability
% Insert in "Experimental Setup" (paragraph 1) and after Table~\ref{tab:sota} (paragraph 2)
% Placeholders [..] are filled from scripts_local/parse_variability.py output.
% ---------------------------------------------------------------

\textbf{Measurement protocol.}
For every (dataset, contraction) pair, GSPARC is executed as five independent
processes. Within each process the contraction is performed three times: the
first execution is a warm-up that absorbs one-time costs (CUDA context creation,
memory-pool initialization, and first-touch page faults) and is excluded from
timing, and the process reports the mean of the remaining two executions. All
GSPARC numbers in the paper are the mean over the five processes; Table~\ref{tab:variability}
reports the corresponding standard deviations.
The CPU baselines (SParTA, SB-TC) and the GPU baseline (GSpTC) do not provide
a warm-up mode; we execute each of them as [three] independent processes and report
the mean of the times measured by their own internal timers, which exclude file
I/O. Every execution is subject to a limit of 3,600 seconds; executions that
exceed it are reported as T.O., and executions that exhaust host or device
memory are reported as OOM. Speedups are computed only over runs that
completed, and the completion rate of each system is reported separately in
Table~\ref{tab:completion}.

\textbf{Run-to-run variability.}
Table~\ref{tab:variability} summarizes the variability of GSPARC across the
five independent runs. The coefficient of variation (CV) of the end-to-end time
is below [X]\% for [N-k] of the [N] cases and at most [Y]\% overall.
The variability is almost entirely attributable to the host-side SLITOM
construction stage, whose multi-threaded sort is sensitive to CPU contention on
the shared server (CV up to [Z]\%), whereas the GPU stages---index matching and
accumulation---are highly reproducible (CV below [W]\% in all cases).
Since the reported speedups over the baselines are one to three orders of
magnitude, this level of variability does not affect any of the conclusions;
we therefore omit error bars from Figures~\ref{fig:sota-frostt}--\ref{fig:sota-ccsd}
for readability and refer to Table~\ref{tab:variability} for the standard deviations.
